\section{Problem Setup}
%Explanations for, for example, the problem formulation, general statement of FTPL, and original GR.
% Maybe other sectioning look better, and it is not mandatory to follow the above line.

In this chapter, we formulate the problem and introduce the framework of FTPL with geometric resampling. At each round $t\in\left[T\right]=\qty{1,2,\dots,T}$, the environment determines a loss vector $\ell_t = (\ell_{t,1}, \ell_{t,2}, \ldots, \ell_{t,K})^\top \in [0,1]^K$. The player then pulls an arm $I_t\in\left[K\right]$, and observes the incurred loss $\ell_{t,I_t}$ associated with the chosen arm.

The loss vector is determined in either a stochastic or adversarial way. In the stochastic setting,
loss vectors $\ell_1, \ell_2, \ldots, \ell_T \in [0,1]^K$ are i.i.d. from an unknown but fixed distribution $\mathcal{P}$ over $\qty[0,1]^K$. The expected loss from arm $i$ is denoted by $\mu_i=\mathbb{E}_{\ell \sim \mathcal{P}}\left[\ell_i\right]\in\left[0,1\right]$. The suboptimality gap of arm $i$ is $\Delta_i=\mu_i-\mu_{i^*}$, where $i^*\in{\arg\min}_{i\in\left[K\right]}\mu_i$ represents the optimal arm.
In this thesis, we assume that $i^*$ is unique in the stochastic setting.
In adversarial setting, the loss $\ell_t$ is determined adversarially, which may depend on the history of loss vectors and chosen arm $\{\left(\ell_s,I_s\right)\}_{s=1}^{t-1}$.

We measure the performance of the player in terms of the pseudo-regret $\mathcal{R}(T)$ defined as
\begin{align*}
    \mathcal{R}(T) = \mathbb{E} \qty[ \sum_{t=1}^T \qty( \ell_{t, I_t} - \ell_{t, i^*} )],\quad
i^* \in \mathop{\arg\min}\limits_{i \in [K]} \mathbb{E} \qty[ \sum_{t=1}^T \ell_{t, i}].
\end{align*}

\subsection{Follow-the-Perturbed-Leader}
We consider the Follow-the-Perturbed-Leader (FTPL) policy,
whose entire procedure is in Algorithm~\ref{alg:FTPL}.
At every round $t$, FTPL keeps the estimated cumulative loss $\hat{L}_t\in\mathbb{R}^K$ specified later
and draws a perturbation vector $r_t=(r_{t,1}, r_{t,2},\dots,r_{t,K})\in\mathbb{R}^K$ from some distribution.
Then, it pulls the arm minimizing the perturbed loss $\hat{L}_{t,i}-r_t/\eta_t$,
where $\eta_t$ is the learning rate.
Unless specified otherwise, we assume that the components of $r_t= (r_{t,1}, r_{t,2}, \cdots, r_{t,K})$ are i.i.d. from
a common distribution $\mathcal{D}$ over $\mathbb{R}$,
whose density function and cumulative distribution function (CDF) are denoted respectively by $f(x)$ and $F(x)$.
In particular,
when $\mathcal{D}$ is supported on $[\nu,\infty)$ for some finite $\nu>-\infty$, we assume $\nu=0$ without loss of generality
by considering the location shift of $\mathcal{D}$.

%At every round $t$, FTPL draws a perturbation $r_t/\eta_t\in\mathbb{R}^K$ to the estimated cumulative loss $\hat{L}_t$ and then pulls the arm that minimizes the perturbed loss. Here $\eta_t = O\qty(t^{-1/2})$ is the exploration parameter or the learning rate. The components of $r_t= (r_{t,1}, r_{t,2}, \cdots, r_{t,K})$ are i.i.d. from a common distribution $\mathcal{D}$, whose density function and cumulative distribution function are denoted respectively as $f(x)$ and $F(x)$.



%However, in certain cases within our work, we employ a biased estimator to enhance computational efficiency. We will show that the introduced bias does not substantially compromise the regret performance of the algorithm.


%\LinesNumbered
%\SetAlgoVlined
\setlength{\textfloatsep}{1.3em}
\begin{algorithm}[t]
\caption{Follow-the-Perturbed-Leader}
\label{alg:FTPL}

\KwIn{Learning rate $\eta_t$}
$\hat{L}_1 \coloneqq 0$\;

\For{$t = 1, 2, \dots, T$}{
    Sample $r_t = (r_{t,1}, r_{t,2}, \dots, r_{t,K})$ i.i.d.\ from $\mathcal{D}$\;
    Pull arm $I_t = {\arg\min}_{i \in [K]} \{\hat{L}_{t,i} - r_{t,i} / \eta_t\}$ and observe $\ell_{t,I_t}$\;
    Compute an estimator $\widehat{w_{t,I_t}^{-1}}$ for $w_{t,I_t}^{-1}$ by geometric resampling\;
    Set $\hat{\ell}_t \coloneqq \ell_{t,I_t} \widehat{w_{t,I_t}^{-1}} e_{I_t}$ and $\hat{L}_{t+1} \coloneqq \hat{L}_t + \hat{\ell}_t$\;
}
\end{algorithm}
% \begin{algorithm}[t]
%     \caption{Follow-the-Perturbed-Leader}
%     \label{alg:FTPL}
%     \begin{algorithmic}[1]
%     \STATE {\bfseries Input:} Learning rate $\eta_t$
%     \STATE Set $\hat{L}_1 \coloneqq 0$
%     \FOR{$t = 1, 2, \dots, T$}
%         \STATE Sample $r_t = (r_{t,1}, r_{t,2}, \dots, r_{t,K})$ i.i.d. from $\mathcal{D}$
%         \STATE Pull arm $I_t = {\arg\min}_{i \in [K]} \{\hat{L}_{t,i} - r_{t,i} / \eta_t\}$ and observe $\ell_{t,I_t}$.
%         \STATE Compute an estimator $\widehat{w_{t,I_t}^{-1}}$ for $w_{t,I_t}^{-1}$ by geometric resampling.
%         \STATE Set $\hat{\ell}_t \coloneqq \ell_{t,I_t} \widehat{w_{t,I_t}^{-1}} e_{I_t}$ and $\hat{L}_{t+1} \coloneqq \hat{L}_t + \hat{\ell}_t$
%     \ENDFOR
%     \end{algorithmic}
% \end{algorithm}




Given $\hat{L}_t$, the probability of pulling arm $i$ is given by
\begin{align}
w_{t,i}
&= \mathbb{P}_{r \sim \mathcal{D}} \left[ i = \mathop{\arg\min}_{j \in [K]} \{\hat{L}_{t,j} - r_j/\eta_t\} \right]\nn
&= \int_{\mathbb{R}} f(z + \eta_t L_{t,i})\prod_{j \neq i} F(z + \eta_t L_{t,j}) \, \mathrm{d}z.
\label{eq_phi}
%&= \int_{\mathbb{R}} f\qty(z + \underline{\lambda}_i)\prod_{j \neq i} F\qty(z + \underline{\lambda}_j) \, \mathrm{d}z,
\end{align}
When $\mathcal{D}$ is supported over $[0,\infty)$
we can write
\begin{align}
w_{t,i}
&= \int_{0}^\infty f(z + \eta_t\underline{L}_{t,i})\prod_{j \neq i} F(z + \eta_t\underline{L}_{t,j}) \, \mathrm{d}z.
\n
%\label{eq_phi}
\end{align}
Here the gap of a vector from its minimum is expressed by underlines, for example, $\underline{L}_{t}=L_t-\mathbf{1} \min_{i\in\left[K\right]} L_{t,i}\in [0, \infty)^K$, where $\mathbf{1}$ is the all-one vector.


%Given $\hat{L}_t$, the probability of pulling arm $i$ is denoted by $w_{t,i}=\phi_i\qty(\eta_t\hat{L}_t)$, where for $\lambda\in\left[0,\infty\right)^K$,
%\begin{align}
%\phi_i\qty(\lambda) &\coloneqq \mathbb{P}_{r \sim \mathcal{D}} \left[ i = \mathop{\arg\min}_{j \in [K]} \{\lambda_j - r_j\} \right] \nonumber \\
%%&= \int_{\nu - \min_{j \in [K]} \lambda_j}^\infty f\qty(z + \lambda_i)\prod_{j \neq i} F\qty(z + \lambda_j) \, \mathrm{d}z \nonumber \\
%&= \int_{\nu}^\infty f\qty(z + \underline{\lambda}_i)\prod_{j \neq i} F\qty(z + \underline{\lambda}_j) \, \mathrm{d}z,
%\label{eq_phi}
%\end{align}
%where $\nu$ refers to the left endpoint of the support of $F$. Here the gap of a vector from its minimum is expressed by underlines, for example, $\underline{\lambda}=\lambda-\mathbf{1} \min_{i\in\left[K\right]} \lambda_i\in [0, \infty)^K$, where $\mathbf{1}$ is the all-one vector.


\subsection{Geometric Resampling}
Since the loss in every round is partially observable,
both FTRL and FTPL generally use an estimator $\hat{\ell}_t$ for loss vector $\ell_t$.
The estimator of the cumulative loss $L_t= \sum_{s=1}^{t-1} \ell_s$ is then obtained as $\hat{L}_t=\sum_{s=1}^{t-1} \hat{\ell}_t$.
%and the entire procedure is given in Algorithm~\ref{alg:FTPL}.
%Denote the cumulative loss at round $t$ by $L_t = \sum_{s=1}^{t-1} \ell_t$, then an estimation of $L_t$ is given by $\hat{L}_t=\sum_{s=1}^{t-1} \hat{\ell}_t$.
%The entire procedure of FTPL is given in Algorithm~\ref{alg:FTPL}.


In many policies for the adversarial setting like FTRL,
the Importance-Weighted (IW) estimator $\hat{\ell}_t=\left( \ell_{t, I_t} / w_{t, I_t}
\right) e_{I_t}$
is employed
as an unbiased estimator for $\ell_t$,
where $e_i$ is a unit vector with the $i$-th component set to one.
% and the others set to zero.
However,
as we can see from \eqref{eq_phi},
$w_{t,I_t}$ is not explicitly computed in FTPL unlike FTRL.
To address this limitation, FTPL instead uses an estimator $\widehat{w_{t,i}^{-1}}$ 
 for $w_{t,i}^{-1}$ by the technique called Geometric Resampling (GR) described in Algorithm~\ref{alg:GR}.
\setlength{\textfloatsep}{1.3em}
\begin{algorithm}[t]
\caption{Geometric Resampling}
\label{alg:GR}

\KwIn{Chosen arm $I_t$, cumulative loss $\hat{L}_t$, learning rate $\eta_t$}
$m \coloneqq 0$\;

\Repeat{$I_t = {\arg\min}_{i \in [K]} \{\hat{L}_{t,i} - r'_{t,i} / \eta_t\}$}{
    $m \coloneqq m + 1$\;
    Sample $r_t' = (r_{t,1}', r_{t,2}', \dots, r_{t,K}')$ i.i.d.\ from $\mathcal{D}$\;
}

$\widehat{w_{t,I_t}^{-1}} \coloneqq m$\;
\end{algorithm}
% \begin{algorithm}[t]
%     \caption{Geometric Resampling}
%     \label{alg:GR}
%     \begin{algorithmic}[1]
%         \STATE \textbf{Input:} \\
%         Chosen arm $I_t$, cumulative loss $\hat{L}_t$, learning rate $\eta_t$
%         \STATE Set $m \coloneqq 0$
%         \REPEAT
%             \STATE $m \coloneqq m + 1$
%             \STATE Sample $r_t^{\prime} = (r_{t,1}^{\prime}, r_{t,2}^{\prime}, \dots, r_{t,K}^{\prime})$ i.i.d. from $\mathcal{D}$
%         \UNTIL{$I_t = {\arg\min}_{i \in [K]} \{\hat{L}_{t,i} - r'_{t,i} / \eta_t\}$}
%         \STATE Set $\widehat{w_{t,I_t}^{-1}} \coloneqq m$
%     \end{algorithmic}
% \end{algorithm}


In GR
employed in \citet{JMLR:v17:15-091} and \citet{pmlr-v201-honda23a},
another perturbation vector $r_t^{\prime}$ is repeatedly drawn from the same
distribution $\mathcal{D}$ until
%the condition
${\arg\min}_{i \in [K]} \{\hat{L}_{t,i} - r'_{t,i}/\eta_t\} = I_t$ is 
satisfied, that is, resampling from $\mathcal{D}$ is repeated until the
arm that would be pulled under the perturbation $r_t^{\prime}$
coincides with the actually pulled arm $I_t$.
% under $r_t$.


%In GR
%%the geometric resampling
%employed in existing work~\citep{JMLR:v17:15-091,pmlr-v201-honda23a},
%we repeatedly draw another perturbation vector $r_t^{\prime}$ from the same
%distribution $\mathcal{D}$ until the condition ${\arg\min}_{i \in [K]} \{\hat{L}_{t,i} - r'_{t,i}/\eta_t\} = I_t$ is 
%satisfied, that is, we repeat the resampling from $\mathcal{D}$ until the
%arm that would be pulled under the perturbation $r_t^{\prime}$
%coincides with the actually pulled arm $I_t$ under $r_t$.


Here the chosen arm $I_t$ is determined prior to the initiation of the GR process. Since the stopping condition is satisfied with probability $w_{t,I_t}$ in each iteration, the number of resampling follows geometric distribution with success probability $w_{t,I_t}$, whose expectation is $1/w_{t,I_t}$.
Letting $\widehat{w_{t,I_t}^{-1}}$ be the number $m$ of resampling,
%Letting $\widehat{w_{t,I_t}^{-1}}\coloneqq m$ for the number $m$ of resampling,
we observe that 
$\widehat{w_{t,I_t}^{-1}}$ serves as an unbiased estimator for $w_{t,I_t}^{-1}$.

Denote the number of resampling taken by geometric resampling at round $t$ as $M_t$,
which is equal to $\widehat{w_{t,I_t}^{-1}}$ in the original GR explained here.
Then, the expectation of $M_t$ given $\hat{L}_t$ is expressed as
\begin{align}
    \mathbb{E}[M_t|\hat{L}_t] &= \sum_{i\in[K]} \mathbb{P}[I_t=i|\hat{L}_t] \mathbb{E}[M_t|\hat{L}_t,I_t=i] \nonumber \\
    &= \sum_{i\in[K]} w_{t,i}\cdot\frac{1}{w_{t,i}}= K.\n
%    &= \sum_{i\in[K]} w_{t,i}\cdot\frac{1}{w_{t,i}} \nonumber \\
%    &= K,
\end{align}
%which is independent of $\hat{L}_t$.
Since generating a perturbation vector $r_t$ from $\mathcal{D}$ needs $K$ random number generations,
the expected complexity per round becomes $O(K^2)$.
%Note that FTPL with geometric resampling generates a perturbation in $O\qty(K)$ time,
%leading to an amortized complexity of $O\qty(K^2)$ per round.
From the viewpoint of the worst-case complexity,
%When we are interested in the worst-case complexity,
it is known that $O(\sqrt{KT\log K})$ regret is achievable
even if we terminate the resampling after $O(\sqrt{KT})$ repetitions, which leads to
$O(\sqrt{K^3T})$
%$O\qty(K^{3/2}T^{1/2})$
worst-case complexity \citep{JMLR:v17:15-091}.
% \honda{please fill values and citations.}
%If we are interested in achieving $O\qty(\sqrt{KT\log K})$ regret, then the process of geometric resampling can be terminated after a finite number of iterations.
On the other hand,
the per-round complexity of FTRL is usually $O(K)$ if the optimization stops within $O(1)$ iterations.
Though FTPL is still efficient for moderate size of $K$ \citep{pmlr-v201-honda23a} thanks to its optimization-free nature,
the worse dependence on $K$ motivates us to develop a more efficient version of FTPL.

